\documentclass[11pt,a4paper]{article}

% ------------------------------------------------------------------ %
% Paquetes
% ------------------------------------------------------------------ %
\usepackage[margin=2.2cm]{geometry}
\usepackage[spanish,es-noquoting]{babel}
\usepackage[utf8]{inputenc}
\usepackage{amsmath,amssymb,mathtools}
\usepackage{booktabs}
\usepackage{longtable}
\usepackage{array}
\usepackage{xcolor}
\usepackage{listings}
\usepackage{siunitx}
\usepackage{microtype}
\usepackage{enumitem}
\usepackage{algorithm}
\usepackage{algpseudocode}
\usepackage[hidelinks]{hyperref}

% ------------------------------------------------------------------ %
% Estilo de listado para fragmentos de parámetros
% ------------------------------------------------------------------ %
\lstdefinestyle{cfg}{
  basicstyle=\ttfamily\footnotesize,
  breaklines=true,
  columns=fullflexible,
  frame=single,
  rulecolor=\color{black!30},
  showstringspaces=false,
}
\lstset{style=cfg}

% Display algorítmico más compacto
\algrenewcommand\algorithmicrequire{\textbf{Entrada:}}
\algrenewcommand\algorithmicensure{\textbf{Salida:}}
\newcommand{\pyconst}[1]{\textcolor{black!60!blue}{\texttt{\footnotesize #1}}}
\newcommand{\dictkey}[1]{\textcolor{black!60!green}{\texttt{\footnotesize #1}}}

% ------------------------------------------------------------------ %
% Macros de conveniencia
% ------------------------------------------------------------------ %
\newcommand{\Vm}{V_{m}}
\newcommand{\Vmsub}[1]{V_{m,#1}}
\newcommand{\Iion}{I_{\text{ion}}}
\newcommand{\Iionsub}[1]{I_{\text{ion},#1}}
\newcommand{\Iext}{I_{\text{ext}}}
\newcommand{\Dten}{\mathbf{D}}
\newcommand{\grad}{\nabla}
\newcommand{\divg}{\nabla\!\cdot}
\newcommand{\Aimp}{\mathsf{A}}
\newcommand{\Bimp}{\mathsf{B}}
\newcommand{\Mmass}{\mathsf{M}}
\newcommand{\Kstif}{\mathsf{K}}
\newcommand{\Lop}{\mathsf{L}}
\newcommand{\code}[1]{\texttt{\small #1}}

\title{Especificación del solver:\\
       \texorpdfstring{\texttt{highOrderElectroActivationFoamImplicitJfNK}}{highOrderElectroActivationFoamImplicitJfNK}\\
       \large Un solver monodominio de alto orden, implícito, Jacobian-free
       Newton--Krylov en OpenFOAM con back-end de álgebra lineal Eigen}
\author{Pablo}
\date{\today}

\begin{document}
\maketitle

\begin{abstract}
Este reporte especifica los algoritmos implementados en el solver de OpenFOAM
\code{highOrderElectroActivationFoamImplicitJfNK}. El solver avanza las ecuaciones
monodominio de la electrofisiología cardíaca acopladas al modelo celular iónico
ten Tusscher--Noble--Noble--Panfilov (TNNP, 2004). El operador de difusión se
discretiza con el marco de volúmenes finitos centrados en celda de OpenFOAM,
opcionalmente mejorado a una reconstrucción de alto orden Local
Radial-Extrapolation (LRE); la fuente iónica se evalúa ya sea en los centros de
celda o en los puntos de cuadratura de celda de la LRE, con el ODE de reacción
integrado en los puntos de cuadratura, una vez por celda y reconstruido, o por un
híbrido consciente del frente (\emph{front-aware}) que hace lo primero solo en el
frente propagante y lo segundo en el resto (Sección~\ref{sec:stateint}); la
integración temporal sigue un esquema $\theta$ (backward Euler o Crank--Nicolson);
el sistema no lineal acoplado resultante puede resolverse por iteración de punto
fijo de Picard, una linealización diagonal de $\Iion$, o un método Jacobian-Free
Newton--Krylov (JFNK) con búsqueda de línea (\emph{line search}) de retroceso de
Armijo, valor inicial por extrapolación lineal y un \emph{clamp} de bracket físico
sobre la entrada al ODE iónico; el sistema lineal interno de JFNK se resuelve con
un GMRES reiniciado propio con ortogonalización de Gram--Schmidt modificada,
construido sobre la biblioteca de álgebra lineal dispersa Eigen; el ODE iónico
por celda usa el integrador adaptativo RKF45 de OpenFOAM con tolerancias
configurables; los tiempos de activación se registran por interpolación lineal a
través del primer cruce de umbral y los puntos del benchmark se muestrean por
ponderación de distancia inversa. Todos los parámetros algorítmicos se exponen a
través de los diccionarios \code{constant/spatialIntegrationProperties} y
\code{constant/timeIntegrationProperties}, y el driver Python compañero
\code{run\_convergence.py} genera estos diccionarios a partir de un único archivo
de configuración. Este solver es algorítmicamente idéntico a su hermano basado en
PETSc \code{highOrderElectroActivationFoamImplicitJfNK\_PETSc}; la única diferencia
es la biblioteca de álgebra lineal (Eigen vs.\ PETSc).
\end{abstract}

\tableofcontents
\bigskip

% ================================================================== %
\section{Notación}
% ================================================================== %

\begin{center}
\begin{tabular}{@{}lll@{}}
\toprule
Símbolo & Significado & Unidades \\
\midrule
$\Omega \subset \mathbb{R}^{d}$ & Dominio espacial ($d \in \{2,3\}$) & \si{\meter} \\
$\partial\Omega$ & Frontera aislada (Neumann de flujo nulo) & --- \\
$t \in (0, T]$ & Tiempo & \si{\second} \\
$\Vm(\mathbf{x}, t)$ & Voltaje transmembrana & \si{\volt} \\
$\mathbf{s}(\mathbf{x}, t) \in \mathbb{R}^{17}$ & States iónicos del TNNP & mixtas \\
$\Iion(\Vm, \mathbf{s})$ & Densidad de corriente iónica total & \si{\volt\per\second} \\
$\Iext(\mathbf{x}, t)$ & Densidad de corriente de estímulo externo & \si{\ampere\per\meter\cubed} \\
$\chi$ & Razón superficie-a-volumen de membrana & \si{\per\meter} \\
$C_{m}$ & Capacitancia de membrana por unidad de área & \si{\farad\per\meter\squared} \\
$\Dten(\mathbf{x})$ & Tensor de conductividad de volumen & \si{\siemens\per\meter} \\
$\theta \in \{0.5, 1\}$ & Parámetro del avance temporal (CN o BE) & --- \\
$\Delta t$ & Longitud del paso de tiempo externo & \si{\second} \\
$N_{c}$ & Número de celdas de malla & --- \\
$\mathsf{V}^{n} \in \mathbb{R}^{N_{c}}$ & Vector de $\Vm$ promediado en celda en $t^{n}$ & \si{\volt} \\
$\Mmass, \Kstif, \Lop$ & Matriz de masa, de rigidez, de difusión escalada & --- \\
$\Aimp, \Bimp$ & Operadores implícitos LHS / RHS del esquema $\theta$ & --- \\
$\mathsf{F}(\mathsf{x})$ & Residuo no lineal discreto a anular & --- \\
\bottomrule
\end{tabular}
\end{center}

% ================================================================== %
\section{Ecuaciones de gobierno}
% ================================================================== %

\subsection{PDE monodominio}

El voltaje de volumen se propaga por la ecuación estándar monodominio de
reacción--difusión~\cite{ttp2004,niederer2012}. En notación física de densidad de
corriente puede escribirse
\begin{equation}
  \chi \, C_{m} \, \frac{\partial \Vm}{\partial t}
   \;=\;
   \divg\!\bigl( \Dten \grad \Vm \bigr)
   - \chi C_m\,\Iion(\Vm, \mathbf{s})
   + \Iext
   \quad \text{en } \Omega \times (0, T],
   \label{eq:monodomain}
\end{equation}
sujeta a fronteras aisladas
$\mathbf{n} \cdot ( \Dten \grad \Vm ) = 0$ en
$\partial \Omega \times (0, T]$, y condición inicial
$\Vm(\mathbf{x}, 0) = -\SI{84}{\milli\volt}$. El tensor de conductividad
$\Dten$ es simétrico definido positivo y ya sea uniforme o leído como un
\code{volTensorField} desde disco. El factor $\chi C_m$ que multiplica a
$\Iion$ en~\eqref{eq:monodomain} refleja la convención usada por la interfaz del
modelo iónico: el solver almacena \code{Iion} como una tasa de voltaje con
dimensiones \si{\volt\per\second}, no como una corriente volumétrica
\si{\ampere\per\meter\cubed}.

Dividiendo~\eqref{eq:monodomain} por $\chi C_{m}$ se obtiene la forma realmente
usada dentro del solver,
\begin{equation}
  \frac{\partial \Vm}{\partial t}
   \;=\;
   \mathcal{L} \Vm
   \;+\;
   s(\Vm, \mathbf{s}),
   \qquad
   \mathcal{L} \;=\; \frac{1}{\chi C_{m}}\,\divg(\Dten \grad \cdot),
   \quad
   s \;=\; -\Iion + \frac{\Iext}{\chi C_{m}}.
   \label{eq:semi-discrete}
\end{equation}
Esto es exactamente lo que \code{rhsFromIion} ensambla en el solver.

\subsection{ODE iónico}

En cada punto de muestreo espacial el vector de states del TNNP
$\mathbf{s} \in \mathbb{R}^{17}$ evoluciona según
\begin{equation}
  \frac{\mathrm{d}\mathbf{s}}{\mathrm{d}t}
   \;=\;
   \mathbf{f}_{\text{TNNP}}\!\bigl(\Vm, \mathbf{s}; \texttt{tissue}\bigr),
   \qquad \mathbf{s}(0) = \mathbf{s}_{0},
   \label{eq:tnnp}
\end{equation}
donde \code{tissue} selecciona \code{epicardialCells}, \code{mCells} o
\code{endocardialCells}. Los 17 states rastreados son
\[
  \mathbf{s} = (V,\ K_{i},\ Na_{i},\ Ca_{i},\ Xr_{1},\ Xr_{2},\ Xs,\ m,\ h,\ j,\ d,\ f,\ fCa,\ s,\ r,\ g,\ Ca_{SR}).
\]
La corriente iónica total devuelta al PDE es la \code{Iion\_cm} algebraica del
TNNP~\cite{ttp2004}.

\subsection{Ecuación MMS manufacturada usada para transferir la convergencia}

El driver de convergencia en
\code{tutorials\_electro/highOrderManufacturedFDAImplicitJfNK} resuelve un problema
manufacturado con la misma estructura numérica que el solver de electro-activación:
PDE implícito tipo monodominio, ODEs locales de states, fuente iónica no lineal,
reconstrucción LRE opcional para $\Vm$ y $\Iion$, y acoplamiento no lineal
seleccionable \code{Picard}/\code{diagonalIion}/\code{JFNK}. Los campos analíticos
son
\begin{align}
  \Vm^{\star}(\mathbf{x},t) &= \sqrt{1+t}\,F(\mathbf{x}),\\
  u_1^{\star}(\mathbf{x},t) &= (1+t)\,G(\mathbf{x})+\Vm^{\star}(\mathbf{x},t),\\
  u_2^{\star}(\mathbf{x},t) &=
      \frac{1}{(1+t)\sqrt{G(\mathbf{x})}},\\
  u_3^{\star}(\mathbf{x},t) &= 0,
\end{align}
con, por ejemplo,
$F=\cos(\pi x)\cos(2\pi y)$ y $G=1+x\,y^2$ en 2D. En 3D,
$F=\cos(\pi x)\cos(2\pi y)\cos(3\pi z)$ y
$G=1+x\,y^2 z^3$. Las ecuaciones de states son
\begin{align}
  \dot u_1 &=
      (u_1+u_3-\Vm)^2 u_2^2
    + \frac{1}{2}(u_1+u_3-\Vm)u_2^2(\Vm-u_3),\\
  \dot u_2 &= -(u_1+u_3-\Vm)u_2^3,\\
  \dot u_3 &= 0.
\end{align}
La fuente iónica manufacturada es
\begin{equation}
  \Iion^{\mathrm{MMS}}(\Vm,\mathbf{u})
  =
  -\frac{1}{2}(u_1+u_3-\Vm)u_2^2(\Vm-u_3)
  + \frac{\beta}{\chi C_m}(\Vm-u_3),
\end{equation}
y la fuente del PDE usada por el solver manufacturado es
$-\Iion^{\mathrm{MMS}}$. El coeficiente $\beta$ es el autovalor espacial del campo
trigonométrico manufacturado bajo el tensor de conductividad constante:
\[
  \beta =
  \begin{cases}
    -\pi^2 D_{xx}, & d=1,\\
    -\pi^2(D_{xx}+4D_{yy}), & d=2,\\
    -\pi^2(D_{xx}+4D_{yy}+9D_{zz}), & d=3.
  \end{cases}
\]
Así, el caso MMS no busca reproducir la fisiología del TNNP; es un problema
controlado con solución exacta conocida que ejercita los mismos operadores
discretos y caminos de solución no lineal usados por este solver.

\subsection{Estímulo externo}

La corriente externa es un escalón ventaneado en tiempo y en región:
\begin{equation}
  \Iext(\mathbf{x}, t) \;=\;
  \begin{cases}
    I_{0} & \mathbf{x} \in \Omega_{\text{stim}}^{(b)},
            \;\; t \in [t_{0}^{(b)},\, t_{0}^{(b)} + \tau],
            \;\; b = 1, \dots, B, \\
    0     & \text{en otro caso},
  \end{cases}
\end{equation}
con $I_{0} = \SI{50000}{\ampere\per\meter\cubed}$ y
$\tau = \SI{2}{\milli\second}$ como valores por defecto del benchmark de Niederer.

% ================================================================== %
\section{Discretización temporal: el esquema $\theta$}
% ================================================================== %

El tiempo se avanza de $t^{n}$ a $t^{n+1} = t^{n} + \Delta t$ por
\begin{equation}
  \frac{\mathsf{V}^{n+1} - \mathsf{V}^{n}}{\Delta t}
   \;=\;
   \theta\, \Lop\, \mathsf{V}^{n+1}
   + (1-\theta)\, \Lop\, \mathsf{V}^{n}
   + \theta\, \mathbf{s}^{n+1}
   + (1-\theta)\, \mathbf{s}^{n}.
   \label{eq:theta}
\end{equation}

\begin{center}
\begin{tabular}{@{}lcl@{}}
\toprule
\code{implicitScheme} & $\theta$ & Esquema \\
\midrule
\code{backwardEuler} & $1$   & Primer orden, L-estable \\
\code{crankNicolson} & $1/2$ & Segundo orden, A-estable \\
\bottomrule
\end{tabular}
\end{center}

Tras multiplicar por $\Delta t^{-1}$ la implementación forma
\begin{equation}
  \bigl[\Mmass/\Delta t \,-\, \theta\, \Lop\bigr]\, \mathsf{V}^{n+1}
   \;=\;
   \bigl[\Mmass/\Delta t \,+\, (1-\theta)\, \Lop\bigr]\, \mathsf{V}^{n}
   \;+\; \theta\, \mathbf{s}^{n+1}
   \;+\; (1-\theta)\, \mathbf{s}^{n},
   \label{eq:implicit-system}
\end{equation}
con los dos operadores Eigen \code{SparseMatrix<scalar,RowMajor>}
\begin{equation}
  \Aimp \;=\; \Mmass/\Delta t \,-\, \theta\, \Lop,
   \qquad
  \Bimp \;=\; \Mmass/\Delta t \,+\, (1-\theta)\, \Lop.
   \label{eq:AB}
\end{equation}
$\Aimp$ y $\Bimp$ se reconstruyen en cada paso de tiempo copiando $\Mmass$,
escalando por $\Delta t^{-1}$ y sumando $\pm \theta \Lop$. Cuando
$\theta = 1$ (BE) la mitad explícita de la difusión se anula y
$\Bimp = \Mmass/\Delta t$. La difusión escalada es
$\Lop = \Kstif / (\chi C_{m})$.

% ================================================================== %
\section{Discretización espacial}
% ================================================================== %

\subsection{Laplaciano FV ortogonal estándar}

\code{assembleStandardOrthogonalStiffnessMatrix} computa el Laplaciano clásico de
5/7 puntos. Por cada cara interna $f$ con área de cara $\mathbf{S}_{f}$,
\begin{equation}
  a_{f} \;=\;
   \frac{ \lvert \mathbf{S}_{f} \rvert\,
          \bigl( \mathbf{n}_{f} \cdot \Dten_{f} \cdot \hat{\mathbf{d}}_{f} \bigr) }
        { \lvert \mathbf{d}_{f} \rvert },
   \quad
   \Dten_{f} = \tfrac{1}{2}\,(\Dten_{\text{own}} + \Dten_{\text{nei}}),
\end{equation}
aportando $-a_{f}/V_{\text{own}}$ a $\Kstif[\text{own},\text{own}]$,
$+a_{f}/V_{\text{own}}$ a $\Kstif[\text{own},\text{nei}]$, y el par transpuesto a
la fila del vecino. Los patches \code{zeroGradient}/\code{empty} aportan flujo
nulo. Los patches \code{fixedValue} aportan solo a la diagonal del dueño (owner).

\subsection{Difusión LRE de alto orden}

Cuando \code{useHighOrder\_Vm = true} y $d>1$,
\code{assembleHighOrderStiffnessMatrix} integra el flujo en una cuadratura de cara
con gradientes reconstruidos:
\begin{equation}
  \Phi_{f} \;=\;
   \sum_{q \in \mathcal{Q}_{f}}
       w_{q}\, \lvert \mathbf{S}_{f} \rvert\,
       \mathbf{n}_{f} \cdot \Dten_{f} \cdot
       \Bigl(
         \sum_{j \in \mathcal{S}_{f}}
            \mathbf{g}_{f,q,j}\, \mathsf{V}_{j}
       \Bigr),
   \label{eq:HO-flux}
\end{equation}
donde $\mathcal{S}_{f}$ es el stencil de cara de la LRE ($\sim 15$--$50$ celdas)
y $\mathbf{g}_{f,q,j} = $~\code{QRGradFaceGPCoeffs}. La contribución de cada
$\mathsf{V}_{j}$ se suma a la fila \code{own} con peso $+\Phi_{f,j}/V_{\text{own}}$
y a la fila \code{nei} con peso $-\Phi_{f,j}/V_{\text{nei}}$, dando un Laplaciano
semi-definido negativo.

\subsection{Matrices de masa}

Dos opciones:
\paragraph{Lumped (\code{massMatrix lumped}).}
\code{assembleDiagonalMassMatrix} escribe $\Mmass = I$.

\paragraph{Consistente HO (\code{massMatrix consistent}).}
\code{assembleConsistentMassMatrixHO} llena
\begin{equation}
  \Mmass[i, j] \;=\;
   \sum_{q \in \mathcal{Q}_{i}}
       \frac{w_{q}}{\sum_{q'} w_{q'}}\,
       \phi_{j}(\mathbf{x}_{q}),
   \quad
   \phi_{j}(\mathbf{x}_{q})
    = \delta_{ij}
    + \mathbf{g}_{i,j} \cdot \mathbf{d}_{q}
    + \tfrac{1}{2}\, \mathbf{H}_{i,j} \!:\! (\mathbf{d}_{q} \otimes \mathbf{d}_{q})
    + \tfrac{1}{6}\, \mathbf{T}^{(3)}_{i,j} \!\vdots\!
        (\mathbf{d}_{q} \otimes \mathbf{d}_{q} \otimes \mathbf{d}_{q}),
   \label{eq:HO-mass}
\end{equation}
con $\mathbf{d}_{q} = \mathbf{x}_{q} - \mathbf{C}_{i}$. Una salvaguarda degrada
silenciosamente a la forma lumped cuando \code{useHighOrder\_Vm = false} o
$d \le 1$.

\subsection{Reconstrucción de alto orden en los puntos de cuadratura iónicos}

Si \code{useHighOrder\_Iion = true}, el ODE iónico se integra en los puntos de
cuadratura de celda. Vm se reconstruye ahí por la misma expansión de Taylor que
en~\eqref{eq:HO-mass}:
\begin{equation}
  \Vm(\mathbf{x}_{q})
   \;\approx\;
   \Vmsub{c} + (\grad \Vm)_{c} \cdot \mathbf{d}_{q}
   + \tfrac{1}{2}\, \mathbf{H}_{c} \!:\! (\mathbf{d}_{q} \otimes \mathbf{d}_{q})
   + \tfrac{1}{6}\, \mathbf{T}^{(3)}_{c} \!\vdots\!
       (\mathbf{d}_{q} \otimes \mathbf{d}_{q} \otimes \mathbf{d}_{q}),
   \label{eq:HO-vm}
\end{equation}
implementado en \code{reconstructVmAtIionIntegrationPoints}.

\subsection{Promediado de alto orden de $\Iion$ a celda}

\code{averageIionIntegrationPointsToCells} contrae el $\Iion$ por punto de
cuadratura de vuelta a celdas:
\begin{equation}
  \Iionsub{i} \;=\;
   \frac{\sum_{q \in \mathcal{Q}_{i}} w_{q}\, \Iionsub{q}}
        {\sum_{q \in \mathcal{Q}_{i}} w_{q}}.
\end{equation}

% ================================================================== %
\section{Integración de states en los puntos de cuadratura iónicos}
\label{sec:stateint}
% ================================================================== %

La cuadratura de alto orden~\eqref{eq:HO-vm} evalúa $\Iion(\Vm,\mathbf{s})$ en los
puntos de cuadratura de celda $\mathbf{x}_q$, así que también necesita ahí el
vector de states iónicos $\mathbf{s}$. \emph{Cómo} llegan los states a los puntos
de cuadratura lo gobierna \dictkey{stateIntegrationMode} y, ortogonalmente,
\dictkey{frontAwareHybrid}. Esta elección resulta decisiva para un frente cardíaco
propagante, y esta sección documenta tanto el algoritmo como la razón por la que
importa.

% ------------------------------------------------------------------ %
\subsection{Dos modos}
% ------------------------------------------------------------------ %

Sea $\mathcal{Q}_i=\{\mathbf{x}_q\}$ el conjunto de cuadratura de la celda
$\Omega_i$ con pesos $w_q$, y sea $\mathbf{s}_c$ un state centrado en celda.

\paragraph{\code{gaussPointODE} (heredado).}
El ODE rígido~\eqref{eq:tnnp} se integra \emph{independientemente en cada punto de
cuadratura}, manejado por el potencial reconstruido localmente:
\begin{equation}
  \mathbf{s}_q^{\,n+1}
   = \mathrm{ODEsolve}\!\bigl(\mathbf{s}_q^{\,n},\,\Vm(\mathbf{x}_q);\Delta t\bigr),
  \qquad
  \Iionsub{q}=\phi\bigl(\Vm(\mathbf{x}_q),\mathbf{s}_q^{\,n+1}\bigr).
  \label{eq:gpode}
\end{equation}
Costo: $N_q$ integraciones rígidas por celda.

\paragraph{\code{cellCentredReconstruct}.}
El ODE se integra \emph{una vez por celda} en el centro de celda, y los states
centrados en celda resultantes se reconstruyen a los puntos de cuadratura por una
LRE dedicada (\code{LREInterp\_states}, orden fijado por
\dictkey{highOrderCoeffs/LRECoeffs\_states}):
\begin{equation}
  \mathbf{s}_c^{\,n+1}
   = \mathrm{ODEsolve}\!\bigl(\mathbf{s}_c^{\,n},\,\overline{V}_{m,c};\Delta t\bigr),
  \quad
  \widehat{\mathbf{s}}(\mathbf{x}_q)
   = \mathcal{R}^{(p)}\!\bigl[\{\mathbf{s}_{c'}^{\,n+1}\}\bigr](\mathbf{x}_q),
  \quad
  \Iionsub{q}=\phi\bigl(\Vm(\mathbf{x}_q),\widehat{\mathbf{s}}(\mathbf{x}_q)\bigr).
  \label{eq:ccrec}
\end{equation}
Costo: una integración rígida por celda --- un factor $N_q$ más barato, lo cual
para las costosas integraciones TNNP es la motivación. Los dos modos difieren solo
en el argumento del mapa no lineal $\Vm\!\mapsto\!\mathbf{s}$: puntual
$\Vm(\mathbf{x}_q)$ en~\eqref{eq:gpode}, versus el promedio de celda
$\overline{V}_{m,c}$ en~\eqref{eq:ccrec}.

% ------------------------------------------------------------------ %
\subsection{Por qué el ODE de reacción debe resolverse en los puntos de Gauss del frente}
\label{sec:jensen}
% ------------------------------------------------------------------ %

Para un campo de reacción suave los dos modos coinciden hasta el orden de
reconstrucción, y \code{cellCentredReconstruct} es preferible por su velocidad. En
un \emph{frente cardíaco propagante} no coinciden, y solo la evaluación puntual
\eqref{eq:gpode} es admisible. La razón es la fuerte no linealidad de la activación
de sodio.

El upstroke (fase de despolarización rápida) lo lleva
$I_{\mathrm{Na}}=g_{\mathrm{Na}}\,m^{3}h\,j\,(\Vm-E_{\mathrm{Na}})$, cuya compuerta
de activación relaja hacia el sigmoide empinado
\begin{equation}
  m_\infty(\Vm)=\Bigl[1+\exp\!\bigl((V_{1/2}-\Vm)/k\bigr)\Bigr]^{-1},
  \qquad V_{1/2}\approx\SI{-40}{mV},\; k\approx\SI{6}{mV}.
  \label{eq:minf}
\end{equation}
Por debajo de su inflexión ($\Vm<V_{1/2}$) --- exactamente el ``pie'' sub-umbral
donde se deciden la ignición y la conducción --- $m_\infty$ es \emph{convexa}.
Escribiendo $\langle\cdot\rangle_i=\bigl(\sum_q w_q\,\cdot\bigr)/\!\sum_q w_q$ para
el promedio de cuadratura de celda, de modo que $\overline{V}_{m,i}=\langle\Vm\rangle_i$,
los dos modos evalúan la compuerta como
\begin{equation}
  \underbrace{\bigl\langle m_\infty(\Vm)\bigr\rangle_i}_{\text{\code{gaussPointODE}}}
  \;\;\ge\;\;
  \underbrace{m_\infty\!\bigl(\langle\Vm\rangle_i\bigr)}_{\text{\code{cellCentredReconstruct}}},
  \label{eq:jensen}
\end{equation}
por la desigualdad de Jensen para la $m_\infty$ convexa. Como
$I_{\mathrm{Na}}\propto m^{3}$, la brecha se amplifica al cubo: el camino de
promedio de celda \emph{sub-estima sistemáticamente} la corriente entrante
regenerativa justo donde $\Vm$ varía dentro de una celda, es decir en el frente de
onda y el borde del estímulo. Reconstruir los states después del ODE (como
en~\eqref{eq:ccrec}) no puede repararlo --- el sesgo se comete \emph{dentro} del
ODE, cuya entrada ya era el potencial promediado.

\paragraph{Una estimación cuantitativa.}
Sea $\Delta V=\max_q\Vm(\mathbf{x}_q)-\min_q\Vm(\mathbf{x}_q)$ el rango sub-celda
de $\Vm$ a través de una celda que abarca el pie del upstroke, y tratemos a $\Vm$
como aproximadamente uniforme sobre la celda (varianza $\sigma^2\simeq\Delta V^2/12$).
Una expansión de punto medio del promedio da
\begin{equation}
  \bigl\langle m_\infty\bigr\rangle
  - m_\infty(\overline V)
  \;\simeq\; \tfrac12\,m_\infty''(\overline V)\,\sigma^2
  \;=\; m_\infty''(\overline V)\,\frac{\Delta V^{2}}{24}.
\end{equation}
En el pie exponencial $m_\infty\!\approx\!e^{(\Vm-V_{1/2})/k}$, de modo que
$m_\infty''/m_\infty=1/k^{2}$ y, usando $I_{\mathrm{Na}}\propto m^{3}$,
\begin{equation}
  \boxed{\;\frac{\delta I_{\mathrm{Na}}}{I_{\mathrm{Na}}}
  \;\approx\; 3\cdot\frac{\Delta V^{2}}{24\,k^{2}}
  \;=\; \frac{\Delta V^{2}}{8\,k^{2}}\; }.
  \label{eq:gap}
\end{equation}
Con $k=\SI{6}{mV}$ y un rango sub-celda $\Delta V\sim\SI{30}{mV}$ a través de un
frente sub-resuelto, \eqref{eq:gap} da
$\delta I_{\mathrm{Na}}/I_{\mathrm{Na}}\approx 900/288\approx 3$ --- la estimación
linealizada ya es $O(1)$, es decir la evaluación de promedio de celda sub-dispara
la corriente de sodio por un factor de orden unidad o más. Siguen dos modos de
falla, ambos observados y documentados en la nota compañera
\code{stateIntegrationMode\_conductionDiagnosis.tex}:
\begin{enumerate}[topsep=2pt,itemsep=1pt]
  \item \textbf{Falla de ignición:} un estímulo nominal queda efectivamente
    sub-umbral; $\Vm$ sube pero el upstroke regenerativo nunca dispara.
  \item \textbf{Bloqueo de conducción:} una vez encendido, cada celda del frente
    aporta muy poca carga entrante para llevar a sus vecinas aguas abajo al umbral
    --- un desajuste source--sink (fuente--sumidero) inducido numéricamente que
    estanca la onda.
\end{enumerate}
Refinar $\Delta t$ hace el frente más filoso (mayor $\Delta V$) y por lo tanto
\emph{agrava} el déficit, razón por la cual el camino \code{cellCentredReconstruct}
puro falla a $\Delta t$ chico mientras que un sobre-disparo de $\Delta t$ grueso
puede enmascararlo. La conclusión es inevitable: en alto orden sobre un frente
filoso el ODE de reacción \emph{debe} resolverse en los puntos de Gauss de las
celdas del frente; solo los volúmenes restantes (suaves) pueden reconstruirse de
forma barata.

% ------------------------------------------------------------------ %
\subsection{Híbrido consciente del frente (front-aware)}
\label{sec:hybrid}
% ------------------------------------------------------------------ %

\dictkey{frontAwareHybrid} (por defecto \code{false}, ortogonal a
\dictkey{stateIntegrationMode}) realiza exactamente esa conclusión: resuelve el ODE
en los puntos de Gauss en la banda delgada de celdas de \emph{frente} y en el
centro de celda (con reconstrucción) en el resto. Una celda $i$ es de \emph{frente}
cuando su rango sub-celda de $\Vm$ excede \dictkey{hybridFrontVmRange} o está
estimulada,
\begin{equation}
  \text{frente}_i \;\Longleftrightarrow\;
  \Bigl(\max_q\Vm(\mathbf{x}_q)-\min_q\Vm(\mathbf{x}_q)\Bigr)
     > \dictkey{hybridFrontVmRange}
  \;\;\vee\;\; (\Iext)_i\neq 0 .
  \label{eq:frontflag}
\end{equation}
Los states persistentes se llevan en los puntos de Gauss (como en
\code{gaussPointODE}); la reconstrucción de celdas suaves deriva sus valores de
celda promediando por cuadratura ese portador. Por evaluación no lineal la rama es:
\begin{enumerate}[topsep=2pt,itemsep=1pt,label=(\alph*)]
  \item celdas de frente: un ODE por punto de Gauss, $\Vm$ local (como~\eqref{eq:gpode});
  \item celdas suaves: un ODE en el centro de celda desde el state promedio de celda;
  \item reconstruir los states de las celdas suaves a sus puntos de Gauss, dejando
        intactos los states de Gauss del frente;
  \item evaluar $\Iion$ en todos los puntos de Gauss y promediar a la celda.
\end{enumerate}

\paragraph{Congelar por paso de tiempo (evitando un ciclo límite).}
La clasificación de frente~\eqref{eq:frontflag} es un \emph{switch discreto
dependiente de la solución}. Si se recalcula dentro del bucle no lineal, las celdas
cerca del umbral alternan frente$\leftrightarrow$suave entre iteraciones; como las
dos ramas computan $\Iion$ por caminos distintos, el término fuente salta
discretamente y el mapa de punto fijo de Picard se vuelve discontinuo, produciendo
un \emph{ciclo límite de período 2} no convergente (observado como
$\text{iter}_{k+2}\equiv\text{iter}_k$ en \code{nonlinearResiduals.dat}). La
clasificación se computa por lo tanto \emph{una vez por paso de tiempo} desde el
$\Vm^{n}$ committeado y se \emph{congela} para todo el solve no lineal, y la banda
se dilata \dictkey{hybridFrontDilation} capas de celdas (vía
\code{mesh.cellCells()}) para que el frente permanezca dentro de la zona resuelta
en Gauss a medida que avanza. Congelar es seguro porque el frente se mueve mucho
menos que una celda por paso (para el benchmark 2D, $\sim\!\SI{0.09}{mm/ms}$, es
decir $\sim\!0.009$ celda por paso a $\Delta t=\SI{e-4}{s}$, $\Delta x=\SI{0.5}{mm}$).

% ------------------------------------------------------------------ %
\subsection{Limitador de la reconstrucción de states}
\label{sec:limiter}
% ------------------------------------------------------------------ %

La reconstrucción de alto orden ($p3$) de los states rígidos del TNNP a través de
un frente sub-resuelto sobre-dispara (Gibbs), y un state llevado muy fuera de su
rango físico puede desbordar (overflow) el álgebra iónica (observado como un
\code{SIGFPE} de PETSc). \dictkey{stateReconstructLimiter} (por defecto
\code{true}, activo para \code{cellCentredReconstruct} y las celdas suaves del
híbrido, independiente de \dictkey{frontAwareHybrid}) aplica una cota tipo
Barth--Jespersen: cada valor reconstruido se recorta al rango de la celda y sus
vecinas de cara,
\begin{equation}
  \widehat{s}(\mathbf{x}_q)
  \;\leftarrow\;
  \mathrm{clip}\!\Bigl(\widehat{s}(\mathbf{x}_q),\;
     \min_{c'\in\{i\}\cup\mathcal{N}(i)} s_{c'},\;
     \max_{c'\in\{i\}\cup\mathcal{N}(i)} s_{c'}\Bigr).
  \label{eq:limiter}
\end{equation}
No se crea ningún extremo nuevo, así que los states reconstruidos permanecen dentro
del rango físicamente realizado (evolucionado por el ODE) y el álgebra iónica nunca
ve una entrada patológica; en regiones suaves la reconstrucción ya cae dentro del
rango y el limitador queda inactivo, preservando el orden ahí.

% ------------------------------------------------------------------ %
\subsection{Claves de diccionario}
% ------------------------------------------------------------------ %

\begin{center}
\begin{tabular}{@{}lll@{}}
\toprule
Clave (en \code{implicitHOCoeffs}) & Significado & Por defecto \\
\midrule
\code{stateIntegrationMode} & \code{gaussPointODE} / \code{cellCentredReconstruct} & \code{cellCentredReconstruct} \\
\code{frontAwareHybrid} & ODEs en Gauss en el frente, reconstrucción en el resto & \code{false} \\
\code{hybridFrontVmRange} & rango sub-celda de $\Vm$ que marca celda de frente [V] & $5\times10^{-3}$ \\
\code{hybridFrontDilation} & capas de celda agregadas alrededor de la banda de frente & $2$ \\
\code{stateReconstructLimiter} & limitador Barth--Jespersen en la reconstrucción & \code{true} \\
\bottomrule
\end{tabular}
\end{center}
\code{LRECoeffs\_states} (en \code{highOrderCoeffs}) fija el orden de
reconstrucción; es requerido por \code{cellCentredReconstruct} y por el híbrido.

% ================================================================== %
\section{Solver no lineal}
% ================================================================== %

En cada paso de tiempo externo la incógnita es el nuevo vector de voltaje
$\mathsf{x}=\mathsf{V}^{n+1}$. Para cualquier candidato $\mathsf{x}$ el solver
primero reconstruye el campo $\Vm$ correspondiente, luego llama a
\code{evaluateNonlinearFields}. Esa rutina resetea los states iónicos a sus valores
al inicio del paso $\mathbf{s}^{n}$, integra los ODEs locales de $t^n$ a $t^{n+1}$
usando el voltaje candidato, obtiene $\Iion(\mathsf{x})$, y opcionalmente promedia
los valores de puntos de cuadratura de vuelta a celdas. Esto da el residuo no lineal
\begin{equation}
  \mathsf{F}(\mathsf{x})
   \;=\;
   \Aimp\, \mathsf{x}
   \;-\;
   \mathsf{r}\!\bigl( \Iion(\mathsf{x}) \bigr),
   \qquad
   \mathsf{x} = \mathsf{V}^{n+1},
   \label{eq:Fdef}
\end{equation}
con
\begin{equation}
  \mathsf{r}(\Iion)
  =
  \Bimp\,\mathsf{V}^{n}
  + \theta\,\Bigl[-\Iion(\mathsf{x})+\Iext^{n+1}/(\chi C_m)\Bigr]
  + (1-\theta)\,\Bigl[-\Iion^{n}+\Iext^{n}/(\chi C_m)\Bigr].
  \label{eq:rhsIion}
\end{equation}
Los tres métodos no lineales difieren solo en cómo llevan a
$\mathsf{F}(\mathsf{x})$ hacia cero.

\subsection{Residuos compartidos y tests de parada}

Cada método registra los mismos diagnósticos:
\begin{align}
  R_{\mathrm{cpl}} &=
    \frac{\|\mathsf{F}(\mathsf{x}_{k})\|_2}
         {\|\mathsf{x}_{k}\|_2+\epsilon},\\
  R_{\Vm} &=
    \frac{\|\mathsf{V}_{k}-\mathsf{V}_{k-1}\|_2}
         {\|\mathsf{V}_{k-1}\|_2+\epsilon},\\
  R_{\Iion} &=
    \frac{\|I_{\mathrm{ion},k}-I_{\mathrm{ion},k-1}\|_2}
         {\|I_{\mathrm{ion},k-1}\|_2+\epsilon},\\
  R_{\mathbf{s}} &=
    \max_j
    \frac{\|\mathbf{s}_{j,k}-\mathbf{s}_{j,k-1}\|_2}
         {\|\mathbf{s}_{j,k-1}\|_2+\epsilon}.
\end{align}
La implementación real usa \code{relativeL2Difference} para $R_{\Vm}$ y $R_{\Iion}$,
y \code{maxStateRelativeL2Difference} para el máximo de states. Una iteración no
lineal solo puede parar tras \code{minNonlinearIterations}; la convergencia luego
requiere \code{nonlinearTolerance}, \code{nonlinearVmTolerance} y
\code{nonlinearIionTolerance}. El criterio de states
\code{nonlinearStatesTolerance} se incluye solo cuando
\code{nonlinearRequireStatesConvergence=true}.

\subsection{Método de punto fijo de Picard}

La rama Picard trata la fuente como retrasada (lagged) durante el solve del PDE. En
la iteración $k$ los states y la fuente iónica se computan desde $\mathsf{x}_k$,
luego el PDE lineal
\begin{equation}
  \Aimp\,\mathsf{V}_{*}
  =
  \mathsf{r}\!\bigl(\Iion(\mathsf{x}_{k})\bigr)
\end{equation}
se resuelve con el solver disperso Eigen seleccionado por \code{linearSolver}. El
siguiente candidato se relaja,
\begin{equation}
  \mathsf{x}_{k+1}
  =
  \mathsf{x}_{k}
  + \omega(\mathsf{V}_{*}-\mathsf{x}_{k}),
  \qquad
  \omega=\code{nonlinearRelaxation}.
\end{equation}
Tras la actualización el solver re-evalúa states y $\Iion$ en $\mathsf{x}_{k+1}$ y
chequea los residuos de arriba. En este solver de electro-activación, a diferencia
de un estudio Picard MMS puro, un solve Picard fallido hace rollback (reversión) a
menos que \code{nonlinearAcceptUnconverged=true}.

\paragraph{Ventajas.}
Picard es simple, barato por iteración no lineal, y usa el mismo solver lineal del
PDE que la formulación implícita clásica. Es una línea base útil y a menudo es
adecuado cuando $\Delta t$ es chico o la fuente iónica es débilmente no lineal en un
paso de tiempo.

\paragraph{Desventajas.}
La derivada de la fuente respecto a $\Vm$ se ignora en la matriz del PDE, así que la
convergencia puede ser lenta cerca del upstroke del potencial de acción o para
$\Delta t$ grande. Un residuo chico puede requerir muchas iteraciones, y puede
necesitarse sub-relajación para robustez.

\subsection{Linealización diagonal de \texorpdfstring{$\Iion$}{Iion}}

La rama \code{diagonalIion} es Picard más una aproximación local por diferencias
finitas de la derivada de la fuente. Definir
\[
  S(\Vm) = -\Iion(\Vm,\mathbf{s}(\Vm)) + \Iext/(\chi C_m).
\]
El solver aproxima, celda a celda,
\begin{equation}
  S'_i(V_{m,k})
  \;\approx\;
  -\frac{
      I_{\mathrm{ion},i}(V_{m,k}+\varepsilon_d)
      - I_{\mathrm{ion},i}(V_{m,k})
    }{\varepsilon_d},
  \qquad
  \varepsilon_d=\code{diagonalIionEpsilon}.
\end{equation}
El residuo linealizado $\Aimp\mathsf{x}-\theta S(\mathsf{x})$ da el sistema
\begin{equation}
  \left(\Aimp-\theta\,\mathrm{diag}(S'_k)\right)\mathsf{x}_{k+1}
  =
  \mathsf{r}\!\bigl(\Iion(\mathsf{x}_{k})\bigr)
  - \theta\,\mathrm{diag}(S'_k)\,\mathsf{x}_{k}.
  \label{eq:diagonalIion}
\end{equation}
La actualización aún puede relajarse con \code{nonlinearRelaxation}, y se usa la
misma lógica de convergencia y rollback que en Picard.

\paragraph{Ventajas.}
Captura la rigidez local dominante de $\Iion(\Vm)$ sin construir un Jacobiano
completo. A menudo reduce el número de iteraciones no lineales respecto a Picard
preservando una matriz escalar dispersa del PDE.

\paragraph{Desventajas.}
Solo se incluye la sensibilidad diagonal local de celda. Los acoplamientos a través
de los states del ODE, el promediado de cuadratura de alto orden y la
reconstrucción espacial no se representan en la matriz. El paso de diferencias
finitas \code{diagonalIionEpsilon} debe elegirse por encima de los niveles de
ODE/ruido pero suficientemente chico para ser una derivada útil.

\subsection{Jacobian-Free Newton--Krylov}

JFNK aplica Newton a~\eqref{eq:Fdef}:
\[
  \mathsf{J}(\mathsf{x}_k)\boldsymbol{\delta}_k
  =
  -\mathsf{F}(\mathsf{x}_k),
  \qquad
  \mathsf{x}_{k+1}=\mathsf{x}_k+\alpha_k\boldsymbol{\delta}_k,
\]
pero nunca ensambla $\mathsf{J}$. GMRES ve solo productos matriz-vector aproximados
por
\begin{equation}
  \mathsf{J}(\mathsf{x}_{k})\, \mathbf{v}
   \;\approx\;
   \frac{\mathsf{F}(\mathsf{x}_{k} + \varepsilon\, \mathbf{v})
         - \mathsf{F}(\mathsf{x}_{k})}{\varepsilon},
   \qquad
   \varepsilon \;=\; \varepsilon_{0}\,
                     \frac{1 + \lVert \mathsf{x}_{k} \rVert_{2}}
                          {\lVert \mathbf{v} \rVert_{2}},
   \label{eq:fdjac}
\end{equation}
con $\varepsilon_{0}=\code{jfnkEpsilon}$. Cada evaluación de $\mathsf{F}$ vuelve a
correr los ODEs y recomputa $\Iion$ desde el mismo snapshot de state al inicio del
paso. Esto es costoso pero da un residuo no lineal genuinamente acoplado.

\subsection{GMRES reiniciado propio}

\code{solveGMRES} implementa GMRES(m) con Gram--Schmidt modificado. Resuelve el
problema de mínimos cuadrados de Krylov para la corrección de Newton usando a lo
sumo \code{jfnkMaxKrylovIterations} vectores base por reinicio y a lo sumo
\code{jfnkMaxRestarts} reinicios. El test de parada lo controla
\code{jfnkLinearTolerance}. El método es ``matrix-free'' (sin matriz) aquí porque el
operador lineal es la envoltura (shell) de diferencias finitas de~\eqref{eq:fdjac},
no una matriz dispersa almacenada.

\paragraph{Ventajas.}
JFNK ve la dependencia no lineal completa
$\Vm\mapsto\mathbf{s}(\Vm)\mapsto\Iion(\Vm,\mathbf{s})$ y usualmente tiene una
convergencia no lineal mucho mejor que Picard en regímenes rígidos. También evita
formar explícitamente un Jacobiano denso o complicado.

\paragraph{Desventajas.}
Cada vector de Krylov requiere una evaluación adicional del residuo, y por ende una
nueva integración de ODE en cada celda o punto de cuadratura. El método es sensible
a la perturbación de diferencias finitas, a las tolerancias del ODE y a las
elecciones de la búsqueda de línea. En esta implementación Eigen el GMRES interno de
JFNK no tiene precondicionador.

\subsection{Búsqueda de línea de retroceso de Armijo}

La actualización cruda de Newton puede aumentar el residuo no lineal cuando el
modelo lineal es pobre. Si \code{jfnkLineSearch=true}, el solver prueba
\[
  \mathsf{x}_{k+1}=\mathsf{x}_{k}+\alpha\boldsymbol{\delta}_{k},
  \qquad
  \alpha=\code{nonlinearRelaxation},\;
  \alpha/2,\;\alpha/4,\dots
\]
hasta que se cumple la condición de descenso suficiente de Armijo:
\begin{equation}
  \lVert \mathsf{F}(\mathsf{x}_{k}+\alpha\boldsymbol{\delta}_{k})
  \rVert_2^2
  \le
  (1-2c\alpha)\,\lVert \mathsf{F}(\mathsf{x}_{k})\rVert_2^2,
  \qquad c=\code{jfnkArmijoC}.
  \label{eq:armijo}
\end{equation}
El retroceso para tras \code{jfnkLineSearchMaxIter} intentos o cuando la longitud de
paso alcanza \code{jfnkLineSearchAlphaMin}. Esto globaliza Newton: los pasos grandes
de Newton se aceptan cuando reducen el residuo, en caso contrario se amortiguan.

\subsection{Valor inicial, clamp del ODE y rollback}

El primer valor inicial para el paso de tiempo es ya sea $\mathsf{V}^{n}$ o la
extrapolación lineal
\begin{equation}
  \mathsf{x}_{0}
  =
  \mathsf{V}^{n}
  + \frac{\Delta t}{\Delta t_{\mathrm{prev}}}
    \left(\mathsf{V}^{n}-\mathsf{V}^{n-1}\right),
  \label{eq:extrap}
\end{equation}
seleccionado por \code{jfnkInitGuessOrder}. Se recorta a
$[\code{jfnkInitGuessVmMin},\code{jfnkInitGuessVmMax}]$. El mismo bracket también
puede aplicarse a los valores de voltaje pasados al solver del ODE mediante
\code{jfnkClampODEInput}; esto protege al modelo iónico durante las sondas de
diferencias finitas de JFNK.

Si el bucle no lineal alcanza su presupuesto de iteraciones sin satisfacer los tests
de convergencia, \code{nonlinearAcceptUnconverged} decide si el último iterado se
committea. Con el valor por defecto \code{false}, el solver hace rollback de $\Vm$,
$\Iion$ y todos los states iónicos al inicio del paso de tiempo. El rollback evita
que un solve no lineal malo contamine el resto de la simulación de activación.

% ================================================================== %
\section{Solver lineal: back-end Eigen}
% ================================================================== %

El envoltorio \code{solveSparseSystem} maneja cualquiera de dos solvers dispersos de
Eigen sobre el operador explícito $\Aimp$ (caminos Picard / diagonal):

\begin{center}
\begin{tabular}{@{}lll@{}}
\toprule
Clave \code{linearSolver} & Tipo Eigen & Tipo \\
\midrule
\code{SparseLU}  & \code{Eigen::SparseLU} & LU directo \\
\code{BiCGSTAB}  & \code{Eigen::BiCGSTAB} + \code{IncompleteLUT} & Iterativo \\
\bottomrule
\end{tabular}
\end{center}

Ambos comparten las claves de diccionario \code{tolerance} y \code{maxIterations}
(la primera la ignora \code{SparseLU} por ser directo). La rama matrix-free de JFNK
nunca instancia estos envoltorios; en cambio llama al GMRES propio (Sección~5.2)
sobre la envoltura del Jacobiano FD.

% ================================================================== %
\section{Integración del ODE iónico}
% ================================================================== %

Para cada celda o punto de cuadratura de celda, el ODE del TNNP~\eqref{eq:tnnp} se
avanza con el integrador adaptativo de OpenFOAM seleccionado vía la clave
\code{solver}. Por defecto \code{RKF45}~\cite{hairer2}. Claves leídas de
\code{constant/timeIntegrationProperties}:

\begin{center}
\begin{tabular}{@{}lll@{}}
\toprule
Clave & Significado & Por defecto \\
\midrule
\code{solver} & Integrador ODE concreto & \code{RKF45} \\
\code{absTol} & Tolerancia absoluta por state & $10^{-9}$ \\
\code{relTol} & Tolerancia relativa por state & $10^{-6}$ \\
\code{maxSteps} & Tope de sub-pasos por llamada externa & $10^{4}$ \\
\code{initialODEStep} & Sub-paso inicial (ms) & $10^{-5}$ \\
\bottomrule
\end{tabular}
\end{center}

\code{relTol} maneja el piso de ruido del Jacobiano matrix-free:
$\eta_{\mathsf{F}} \sim $~\code{relTol}. Ajustar \code{relTol} más fino compra una
tolerancia externa alcanzable más baja a costa de más sub-pasos de RKF45 por
iteración de Newton.

% ================================================================== %
\section{Protección numérica del TNNP}
% ================================================================== %

El switch maestro \code{TNNPNumericalProtection} (por defecto \code{false})
condiciona una guarda local invocada en cada entrada externa al kernel TNNP embebido
(\code{calculateCurrent}, \code{solveODE\_start}, \code{solveODE\_end},
\code{derivatives}).

\begin{center}
\begin{tabular}{@{}lll@{}}
\toprule
Clave & Tipo & Por defecto \\
\midrule
\code{TNNPNumericalProtection} & Switch & \code{false} \\
\code{TNNPClampGates} & Switch & \code{true} \\
\code{TNNPClampVmForRates} & Switch & \code{true} \\
\code{TNNPLogNumericalProtection} & Switch & \code{true} \\
\code{TNNPConcentrationFloor} & scalar (mM) & $10^{-12}$ \\
\code{TNNPVMinForRates} & scalar (mV) & $-200$ \\
\code{TNNPVMaxForRates} & scalar (mV) & $+100$ \\
\code{TNNPVoltageZeroEps} & scalar (mV) & $10^{-8}$ \\
\code{TNNPMaxProtectionLogEntries} & label & $10^{6}$ \\
\code{TNNPProtectionLog} & word & \code{TNNP\_numericalProtection\_summary.dat} \\
\bottomrule
\end{tabular}
\end{center}

Cuando la protección está habilitada, las guardas aplicadas en orden son:
\begin{enumerate}[topsep=0pt,itemsep=2pt]
  \item \textbf{$V$ no finito}: \code{!std::isfinite(V)} $\rightarrow$
        resetear a $-86.2$ mV.
  \item \textbf{Clamp de $V$ para las tasas} (si \code{TNNPClampVmForRates}):
        recortar a $[\,$\code{TNNPVMinForRates}, \code{TNNPVMaxForRates}$\,]$.
  \item \textbf{Evitación de singularidad en $V=0$}: empujar $V$ fuera de cero por
        \code{TNNPVoltageZeroEps} (preservando el signo) cuando
        $\lvert V \rvert$ cae por debajo de ese umbral.
  \item \textbf{Piso de concentración} en $K_{i}, Na_{i}, Ca_{i}, Ca_{SR}$:
        entradas no finitas o por debajo del piso se resetean a
        \code{TNNPConcentrationFloor}.
  \item \textbf{Clamp de compuertas} en $Xr_{1}, Xr_{2}, Xs, m, h, j, d, f, fCa,
        s, r, g$ (si \code{TNNPClampGates}): no finitas a cero,
        recortar a $[0, 1]$.
\end{enumerate}

La implementación embebida almacena un conteo de correcciones acotado y emite un
archivo de resumen compacto al salir; no depende de utilidades externas de logging
del modelo iónico.

% ================================================================== %
\section{Seguimiento de la activación y benchmarks}
% ================================================================== %

\paragraph{Tiempo de activación por celda.}
El primer cruce de $V_{\text{thr}}$ en un paso de tiempo se registra por
interpolación lineal:
\begin{equation}
  w \;=\; \mathrm{clip}\!\Bigl(
      \frac{V_{\text{thr}} - V^{n}}{V^{n+1} - V^{n}},
      \; 0, 1
    \Bigr),
   \qquad
   t_{\text{act}} \;=\; t^{n} + w\, \Delta t,
\end{equation}
con una bandera por celda que se limpia tras el primer cruce para ignorar cruces
posteriores.

\paragraph{Terminación temprana en P8.}
Cuando \code{stopAfterPointActivation = true}, se monitorea la celda más cercana a
\code{stopActivationPoint} (P8 en el benchmark de Niederer) y \code{effectiveEndTime}
se fija a $t_{\text{act}} + $~\code{stopDelayAfterActivation} una vez que su
$t_{\text{act}} > 0$.

\paragraph{Muestreo de la diagonal.}
\code{sampleIDW} realiza interpolación ponderada por distancia inversa en
\code{nDiagonalSamples} puntos entre \code{diagonalStart} y \code{diagonalEnd}, con
$k = $~\code{sampleNeighbours}, $p = $~\code{samplePower}. Como una muestra de la
diagonal promedia sus $k$ celdas más cercanas y las celdas no activadas contribuyen
$t_{\text{act}}=0$, una muestra observada \emph{en vivo} sube desde abajo a medida
que el frente llena su vecindario; el valor final (las $k$ vecinas activadas) es el
correcto.

\paragraph{Salida de activación en vivo.}
\code{writeActivationSamples} se llama una vez antes del bucle de tiempo (para que
\code{diagonalActivationTime.dat} y \code{P8\_activationTime.dat} existan desde $t=0$
con todo punto pendiente) y de nuevo cada vez que el conteo de celdas activadas
aumenta, sobrescribiendo los archivos. La corrida puede así monitorearse a medida
que el frente avanza; el driver expone el archivo en vivo en cada directorio de
corrida mediante un enlace simbólico que se reemplaza por una copia real al final.

\paragraph{Detección de falla de activación.}
Cuando \dictkey{abortOnActivationFailure = true} el solver para temprano
(limpiamente, exit~0) si la activación falla, evitando una corrida completa hasta
\code{endTime}. Dos criterios, usando el conteo corriente de celdas activadas $N_a$
y el tiempo $t_\star$ de la última activación nueva:
\begin{itemize}[topsep=2pt,itemsep=1pt]
  \item \textbf{sin ignición:} $N_a=0$ y
        $t>\dictkey{activationIgnitionTimeout}$;
  \item \textbf{bloqueo de conducción:} $N_a>0$, el punto de parada aún no se
        alcanzó, y $t-t_\star>\dictkey{activationStallTime}$.
\end{itemize}
El resultado se registra siempre en
\code{postProcessing/.../activationStatus.dat} como uno de \code{OK},
\code{FAILED\_NO\_IGNITION}, \code{FAILED\_CONDUCTION\_BLOCK} o \code{INCOMPLETE}
(con el tiempo final y $N_a$). El driver lee este archivo y, ante una falla, lo
loguea en \code{activation\_status\_summary.dat} y continúa con la siguiente
configuración en lugar de abortar el barrido.

% ================================================================== %
\section{Diagrama de flujo del solver}
% ================================================================== %

El pseudo-código de abajo sigue el solver real. El setup compartido es común a todos
los métodos no lineales; luego se ejecuta uno de los tres bloques específicos según
\dictkey{nonlinearMethod}.

\begin{algorithm}[H]
\caption{Setup compartido del paso de tiempo y commit/rollback final.}
\begin{algorithmic}[1]
\State Leer diccionarios y ensamblar $\Mmass$, $\Kstif$, $\Lop=\Kstif/(\chi C_m)$.
\While{$t^n < T$}
  \State Guardar snapshots: $\mathsf{V}^n$, $\Iion^n$, $\mathbf{s}^n$.
  \State Construir $\Aimp=\Mmass/\Delta t-\theta\Lop$ y
         $\Bimp=\Mmass/\Delta t+(1-\theta)\Lop$.
  \State Construir la fuente vieja $S^n=-\Iion^n+\Iext^n/(\chi C_m)$.
  \If{\dictkey{jfnkInitGuessOrder}$\ge 1$ y existe paso previo}
    \State $\mathsf{x}_0\gets\mathsf{V}^n+
      (\Delta t/\Delta t_{\rm prev})(\mathsf{V}^n-\mathsf{V}^{n-1})$.
    \State Recortar $\mathsf{x}_0$ a
      $[\dictkey{jfnkInitGuessVmMin},\dictkey{jfnkInitGuessVmMax}]$.
  \Else
    \State $\mathsf{x}_0\gets\mathsf{V}^n$.
  \EndIf
  \State Correr el Algoritmo~\ref{alg:picard},~\ref{alg:diagonal} o~\ref{alg:jfnk}.
  \If{no convergió y \dictkey{nonlinearAcceptUnconverged} es false}
    \State Hacer rollback de $\Vm,\Iion,\mathbf{s}$ a los snapshots guardados en $t^n$.
  \Else
    \State Committear el iterado no lineal final.
  \EndIf
  \State Actualizar tiempos de activación y muestras del benchmark.
\EndWhile
\end{algorithmic}
\end{algorithm}

\begin{algorithm}[H]
\caption{Solve no lineal de Picard.}
\label{alg:picard}
\begin{algorithmic}[1]
\Require Candidato inicial $\mathsf{x}_0$, snapshot de state $\mathbf{s}^n$.
\For{$k=0,\dots,\code{nonlinearIterations}-1$}
  \State Guardar iterado previo $\mathsf{x}_{k-1}$, $I_{\mathrm{ion},k-1}$,
         $\mathbf{s}_{k-1}$ para los residuos.
  \State \Call{evaluateNonlinearFields}{$\mathsf{x}_k$}:
    resetear states a $\mathbf{s}^n$, reconstruir $\Vm$ en los puntos de cuadratura
    de Iion si está habilitado, llamar a \code{ionicModel->solveODE}, computar $I_{\mathrm{ion},k}$.
  \State Construir el RHS $\mathsf{r}(I_{\mathrm{ion},k})$ usando~\eqref{eq:rhsIion}.
  \State Resolver $\Aimp\,\mathsf{V}_*=\mathsf{r}(I_{\mathrm{ion},k})$ con
    \dictkey{linearSolver}.
  \State Relajar:
    $\mathsf{x}_{k+1}\gets\mathsf{x}_{k}+
    \code{nonlinearRelaxation}(\mathsf{V}_*-\mathsf{x}_{k})$.
  \State Re-evaluar los ODEs y $\Iion$ en $\mathsf{x}_{k+1}$.
  \State Computar $R_{\mathrm{cpl}}$, $R_{\Vm}$, $R_{\Iion}$ y el opcional
    $R_{\mathbf{s}}$.
  \If{todos los criterios pedidos pasan y $k+1\ge\code{minNonlinearIterations}$}
    \State Marcar convergido y salir del bucle.
  \EndIf
\EndFor
\end{algorithmic}
\end{algorithm}

\begin{algorithm}[H]
\caption{Solve no lineal \code{diagonalIion}.}
\label{alg:diagonal}
\begin{algorithmic}[1]
\Require Candidato inicial $\mathsf{x}_0$, snapshot de state $\mathbf{s}^n$.
\For{$k=0,\dots,\code{nonlinearIterations}-1$}
  \State Evaluar states y $I_{\mathrm{ion},k}$ en $\mathsf{x}_k$ exactamente como en Picard.
  \State Aproximar la derivada local
    $S'_k\approx-[\Iion(\mathsf{x}_k+\varepsilon_d)-\Iion(\mathsf{x}_k)]/\varepsilon_d$,
    con $\varepsilon_d=\dictkey{diagonalIionEpsilon}$.
  \State Ensamblar
    $\Aimp_k=\Aimp-\theta\,\mathrm{diag}(S'_k)$.
  \State Construir el RHS corregido
    $\mathsf{r}_k=\mathsf{r}(I_{\mathrm{ion},k})-\theta\,\mathrm{diag}(S'_k)\mathsf{x}_k$.
  \State Resolver $\Aimp_k\mathsf{V}_*=\mathsf{r}_k$ con \dictkey{linearSolver}.
  \State Relajar, re-evaluar ODEs/$\Iion$, computar residuos y testear convergencia
    como en Picard.
\EndFor
\end{algorithmic}
\end{algorithm}

\begin{algorithm}[H]
\caption{Solve Jacobian-Free Newton--Krylov.}
\label{alg:jfnk}
\begin{algorithmic}[1]
\Require Candidato inicial $\mathsf{x}_0$, snapshot de state $\mathbf{s}^n$.
\For{$k=0,\dots,\code{nonlinearIterations}-1$}
  \State Evaluar $\mathsf{F}_k=\mathsf{F}(\mathsf{x}_k)$; esto resetea states a
    $\mathbf{s}^n$, resuelve los ODEs y recomputa $\Iion$.
  \State Computar residuos. Si todos los criterios pedidos pasan, marcar convergido
    y salir.
  \State Definir el producto matrix-free
    $\mathcal{J}_k\mathbf{v}=
    [\mathsf{F}(\mathsf{x}_k+\varepsilon\mathbf{v})-\mathsf{F}_k]/\varepsilon$,
    con $\varepsilon$ de \dictkey{jfnkEpsilon}.
  \State Resolver $\mathcal{J}_k\boldsymbol{\delta}_k=-\mathsf{F}_k$ por
    GMRES reiniciado usando \dictkey{jfnkMaxKrylovIterations},
    \dictkey{jfnkMaxRestarts} y \dictkey{jfnkLinearTolerance}.
  \If{\dictkey{jfnkLineSearch}}
    \State Comenzar con $\alpha=\code{nonlinearRelaxation}$ y retroceder hasta que
      Armijo~\eqref{eq:armijo} tenga éxito, se agote el presupuesto de intentos,
      o se alcance \dictkey{jfnkLineSearchAlphaMin}.
  \Else
    \State $\alpha\gets\code{nonlinearRelaxation}$.
  \EndIf
  \State $\mathsf{x}_{k+1}\gets\mathsf{x}_{k}+\alpha\boldsymbol{\delta}_k$.
  \State Re-evaluar $\mathsf{F}(\mathsf{x}_{k+1})$, actualizar campos, computar
    residuos y testear convergencia.
\EndFor
\end{algorithmic}
\end{algorithm}

\noindent\textbf{Leyenda.} \pyconst{Azul} = constante Python fijada en
\code{run\_convergence.py}; \dictkey{verde} = clave de diccionario OpenFOAM leída
por el solver. Constantes y claves están mapeadas uno-a-uno en la
sección~\ref{sec:driver}.

% ================================================================== %
\section{Driver manufacturado: \texttt{run\_convergence.py}}
\label{sec:driver}
% ================================================================== %

El script de convergencia
\code{/home/pablo/OpenFOAM/pablo-v2312/run/tutorials\_electro/highOrderManufacturedFDAImplicitJfNK/run\_convergence.py}
genera un caso MMS cuyo camino numérico busca coincidir con este solver. Las tablas
de abajo describen las constantes expuestas ahí, las entradas de diccionario OpenFOAM
que generan, y qué parte de las ecuaciones discretas afectan. Las mismas claves de
diccionario las lee \code{highOrderElectroActivationFoamImplicitJfNK}; solo difieren
ítems específicos del problema como campos exactos versus TNNP/estímulo.

\subsection{Geometría y barrido de parámetros}
\begin{longtable}{@{}p{0.30\linewidth}p{0.28\linewidth}p{0.36\linewidth}@{}}
\toprule
Constante Python & Clave OpenFOAM & Ecuación / mecanismo afectado \\
\midrule \endhead
\code{DIMENSIONS} & --- &
   Selecciona el dominio manufacturado 2D o 3D y los stencils LRE
   correspondientes en \code{P\_LEVELS\_BY\_DIM}. \\
\code{N\_VALUES} & conteos de celdas de \code{blockMeshDict} &
   Refinamiento espacial. Cambia el error de discretización FV/LRE en el
   PDE, la cuadratura de Iion y la reconstrucción de states. \\
\code{DT\_VALUES} & \code{controlDict.deltaT} &
   Paso de tiempo externo del PDE. También fija el intervalo de integración de cada
   solve de ODE local. \\
\code{END\_TIME} & \code{controlDict.endTime} &
   Tiempo final del test MMS transitorio. \\
\code{IMPLICIT\_SCHEMES} &
   \code{spatialIntegrationProperties.implicitScheme} &
   Elige $\theta$ en~\eqref{eq:theta}: \code{backwardEuler}
   ($\theta=1$) o \code{crankNicolson} ($\theta=1/2$). \\
\code{MASS\_MATRIX\_VALUES} &
   \code{spatialIntegrationProperties.massMatrix} &
   Elige $\Mmass$ lumped o consistente-LRE en
   $\Aimp=\Mmass/\Delta t-\theta\Lop$. \\
\code{ACTIVE\_HO\_PAIRS\_Vm\_States} &
   \code{useHighOrder\_Vm} / \code{useHighOrder\_Iion}
   + \code{LRECoeffs\_*} &
   Tupla \code{(nivel-Vm, nivel-Iion/states)}. El nivel de Vm afecta la
   difusión y la masa consistente; el nivel de Iion/states afecta dónde se
   resuelven los ODEs y dónde se integra la fuente no lineal. \\
\code{P\_LEVELS\_BY\_DIM} & \code{highOrderCoeffs} &
   Mapea \code{p1,p2,p3} al orden polinómico LRE \code{N}, objetivo de stencil
   \code{Nn}, exponente del peso radial \code{k}, y
   \code{maxStencilSize}. \\
\bottomrule
\end{longtable}

\subsection{Método no lineal y controles de convergencia}
\begin{longtable}{@{}p{0.35\linewidth}p{0.30\linewidth}p{0.30\linewidth}@{}}
\toprule
Constante Python & Clave OpenFOAM & Ecuación / mecanismo afectado \\
\midrule \endhead
\code{NONLINEAR\_METHOD} & \code{nonlinearMethod} &
   Selecciona la estrategia no lineal para el acoplamiento PDE/fuente:
   \code{Picard}, \code{JFNK}, o \code{diagonalIion}. \\
\code{IMPLICIT\_NONLINEAR\_ITERATIONS} & \code{nonlinearIterations} &
   Número máximo de correcciones no lineales por paso de tiempo del PDE. \\
\code{IMPLICIT\_MIN\_NONLINEAR\_ITERATIONS} &
   \code{minNonlinearIterations} &
   Número mínimo de correcciones antes de que un test de convergencia pueda parar el
   bucle. \\
\code{NONLINEAR\_TOLERANCE} & \code{nonlinearTolerance} &
   Tolerancia del residuo acoplado del PDE $R_{\mathrm{cpl}}$. Usada directamente por
   JFNK y diagonalIion; en el solver electro también se reporta para Picard. \\
\code{NONLINEAR\_VM\_TOLERANCE} & \code{nonlinearVmTolerance} &
   Corrección relativa de la incógnita del PDE $\Vm$. \\
\code{NONLINEAR\_STATES\_TOLERANCE} &
   \code{nonlinearStatesTolerance} &
   Corrección relativa de todos los states del ODE. En el solver MMS estos son
   $u_1,u_2,u_3$; en este solver electro son las variables de state del TNNP. \\
\code{NONLINEAR\_REQUIRE\_STATES\_CONVERGENCE} &
   \code{nonlinearRequireStatesConvergence} &
   Si es true, el residuo de states es parte del criterio de parada. \\
\code{NONLINEAR\_ACCEPT\_UNCONVERGED} &
   \code{nonlinearAcceptUnconverged} &
   Si es false, un solve no lineal electro fallido hace rollback de $\Vm,\Iion$ y
   states a $t^n$. \\
\code{NONLINEAR\_IION\_TOLERANCE} & \code{nonlinearIionTolerance} &
   Corrección relativa de la fuente no lineal $\Iion$. \\
\code{NONLINEAR\_RELAXATION} & \code{nonlinearRelaxation} &
   Factor de mezcla Picard/diagonal o longitud de paso inicial de Newton para
   JFNK/Armijo. \\
\bottomrule
\end{longtable}

\subsection{Controles de JFNK y de diagonal-Iion}
\begin{longtable}{@{}p{0.35\linewidth}p{0.30\linewidth}p{0.30\linewidth}@{}}
\toprule
Constante Python & Clave OpenFOAM & Ecuación / mecanismo afectado \\
\midrule \endhead
\code{JFNK\_MAX\_KRYLOV\_ITERATIONS} & \code{jfnkMaxKrylovIterations} &
   Longitud de reinicio $m$ de GMRES(m), controlando la memoria y la calidad de la
   aproximación de Krylov. \\
\code{JFNK\_MAX\_RESTARTS} & \code{jfnkMaxRestarts} &
   Número de veces que la base de GMRES puede descartarse y reconstruirse. \\
\code{JFNK\_LINEAR\_TOLERANCE} & \code{jfnkLinearTolerance} &
   Tolerancia relativa del solve de la corrección de Newton
   $\mathsf{J}\delta=-\mathsf{F}$. \\
\code{JFNK\_EPSILON} & \code{jfnkEpsilon} &
   Perturbación base para los productos matriz-vector matrix-free del Jacobiano.
   Debe estar por encima del piso de ODE/ruido pero suficientemente chica para
   aproximar una derivada. \\
\code{JFNK\_INIT\_GUESS\_ORDER} & \code{jfnkInitGuessOrder} &
   Elige el $\Vm$ inicial retrasado o linealmente extrapolado para el solve no
   lineal. \\
\code{JFNK\_INIT\_GUESS\_VM\_MIN/MAX} &
   \code{jfnkInitGuessVmMin/Max} &
   Intervalo de clamp para el valor inicial extrapolado y, si está habilitado, las
   entradas del ODE. \\
\code{JFNK\_CLAMP\_ODE\_INPUT} & \code{jfnkClampODEInput} &
   Protege el solve del ODE de sondas de diferencias finitas de JFNK no físicas. \\
\code{JFNK\_LINE\_SEARCH} & \code{jfnkLineSearch} &
   Activa/desactiva el retroceso de Armijo. \\
\code{JFNK\_LINE\_SEARCH\_MAX\_ITER} &
   \code{jfnkLineSearchMaxIter} &
   Número máximo de divisiones por la mitad del paso de Newton. \\
\code{JFNK\_LINE\_SEARCH\_ALPHA\_MIN} &
   \code{jfnkLineSearchAlphaMin} &
   Piso para la longitud de paso de Armijo. \\
\code{JFNK\_ARMIJO\_C} & \code{jfnkArmijoC} &
   Parámetro de descenso suficiente en~\eqref{eq:armijo}. \\
\code{DIAGONAL\_IION\_EPSILON} & \code{diagonalIionEpsilon} &
   Paso de diferencias finitas para la derivada diagonal de la fuente en
   \code{diagonalIion}. \\
\bottomrule
\end{longtable}

\subsection{Solver lineal del PDE}
\begin{longtable}{@{}p{0.30\linewidth}p{0.30\linewidth}p{0.35\linewidth}@{}}
\toprule
Constante Python & Clave OpenFOAM & Ecuación / mecanismo afectado \\
\midrule \endhead
\code{IMPLICIT\_LINEAR\_SOLVER} & \code{linearSolver} &
   Solver disperso Eigen usado por Picard y diagonalIion para
   $\Aimp\mathsf{V}=\mathsf{r}$. JFNK en cambio usa GMRES sobre la envoltura
   Jacobian-free. \\
\code{IMPLICIT\_TOLERANCE} & \code{tolerance} &
   Tolerancia de BiCGSTAB para el solve lineal del PDE. \\
\code{IMPLICIT\_MAX\_ITERATIONS} & \code{maxIterations} &
   Tope de iteraciones de BiCGSTAB. \\
\bottomrule
\end{longtable}

\subsection{Integrador del ODE de states}
\begin{longtable}{@{}p{0.30\linewidth}p{0.30\linewidth}p{0.35\linewidth}@{}}
\toprule
Constante Python & Clave OpenFOAM & Ecuación / mecanismo afectado \\
\midrule \endhead
\code{STATE\_ODE\_SOLVER} & \code{stateODESolver} /
   \code{timeIntegrationProperties.solver} &
   Integrador de ODE local para states MMS o states TNNP. \\
\code{STATE\_ODE\_ABS\_TOL} & \code{stateODEAbsTol} /
   \code{timeIntegrationProperties.absTol} &
   Tolerancia absoluta del solve del ODE. \\
\code{STATE\_ODE\_REL\_TOL} & \code{stateODERelTol} /
   \code{timeIntegrationProperties.relTol} &
   Tolerancia relativa del ODE. Controla el piso de ruido visto por las diferencias
   finitas de JFNK. \\
\code{STATE\_ODE\_INITIAL\_STEP} & \code{stateODEInitialStep} /
   \code{initialODEStep} &
   Primer sub-paso RKF45 intentado dentro de un paso de tiempo del PDE. \\
\code{STATE\_ODE\_MAX\_STEPS} & \code{stateODEMaxSteps} /
   \code{maxSteps} &
   Máximo de sub-pasos de ODE por solve local. \\
\bottomrule
\end{longtable}

\subsection{Parámetros LRE de alto orden por nivel-p}
\code{P\_LEVELS\_BY\_DIM[dim][p]} da, para cada (dim, p):
\code{N} (orden polinómico), \code{Nn} (vecinos del stencil),
\code{k} (radio de búsqueda), \code{maxStencilSize} (tope). El script escribe estos
en un bloque anidado
\code{highOrderCoeffs\{LRECoeffs\_Vm\{\dots\}; LRECoeffs\_Iion\{\dots\}\}}
cuando HO está habilitado para Vm o $\Iion$.

\subsection{Diagnósticos y control de corrida}
\begin{longtable}{@{}p{0.35\linewidth}p{0.30\linewidth}p{0.30\linewidth}@{}}
\toprule
Constante Python & Clave OpenFOAM & Mecanismo afectado \\
\midrule \endhead
\code{WRITE\_NONLINEAR\_RESIDUALS} &
   \code{writeNonlinearResiduals} & Log de residuos on/off. \\
\code{KEEP\_LOGS\_ON\_SUCCESS} & --- &
   Switch solo-Python: conservar logs del solver/blockMesh tras corridas exitosas.
   Los logs de corridas fallidas se conservan siempre. \\
\bottomrule
\end{longtable}

Para el tutorial de electro-activación el driver análogo también expone parámetros de
umbral de activación, parada en P8 y muestreo de la diagonal. No afectan el solve
PDE/ODE en sí; afectan solo la parada y el post-procesamiento de los tiempos de
activación.

% ================================================================== %
\section{Resultados}
\label{sec:results}
% ================================================================== %

Esta sección reúne los resultados de convergencia actuales para los benchmarks 2D y
3D estilo Niederer, corridos con las tres estrategias no lineales expuestas por el
solver (\code{Picard}, \code{diagonalIion}, \code{JFNK}). Todas las corridas usan el
esquema temporal Crank--Nicolson ($\theta = 1/2$), la matriz de masa consistente de
alto orden cuando la LRE está habilitada, el integrador de ODE de states
\code{RKF45}, y un único $\Delta t_{\text{req}} = 10^{-4}$~s solicitado. Se muestrean
seis pares LRE $(\Vm,\Iion)$: \code{NO}/\code{NO}, \code{NO}/\code{p1},
\code{NO}/\code{p3}, \code{p3}/\code{NO}, \code{p3}/\code{p1}, \code{p3}/\code{p3}.
Cada tabla lista, para la métrica de su caption, el valor de cada corrida disponible;
\textsc{NC} marca las combinaciones que aún no se han corrido. Los resultados 2D
viven bajo
\path{highOrderNiedererEtAl2012ImplicitJfNK[_diagonalIion|_Picard]/results/example0/2D/};
los resultados 3D bajo los mismos árboles pero con el sufijo \path{3D/}. El tiempo de
reloj de pared (wall-clock) es el tiempo transcurrido reportado por
\texttt{/usr/bin/time -v}, redondeado a minutos enteros. El RSS pico es el tamaño
máximo del conjunto residente (resident set size) del proceso del solver, expresado
en MiB ($\lfloor\text{kB}/1024\rfloor$).

\subsection{Benchmark 2D estilo Niederer}
\label{sec:results:2D}

El caso 2D sigue la reducción planar de 20$\times$8~mm usada por los reportes
explícito e implícito no-JfNK (una sola celda en $z$, fibras planares), ejercitado a
tamaños de malla $\Delta x\in\{0.5, 0.25, 0.125, 0.1\}$~mm; solo el barrido
\code{Picard} incluye actualmente la malla más fina.

\begin{center}\footnotesize
\begin{longtable}{llrrrr}
\caption{Tiempo de activación P8 2D [ms] a $\Delta t = 10^{-4}$~s. Las columnas son
tamaños de malla $\Delta x$ en mm.}\\
\toprule
Vm & Iion & $0.5$ & $0.25$ & $0.125$ & $0.1$ \\
\midrule\endfirsthead\toprule
Vm & Iion & $0.5$ & $0.25$ & $0.125$ & $0.1$ \\
\midrule\endhead
\multicolumn{6}{l}{\textit{\code{Picard}}}\\
NO & NO & 153.32 & 62.99 & 45.03 & 42.76 \\
NO & p1 &  91.05 & 54.86 & 43.21 & 41.84 \\
NO & p3 &  32.46 & 39.14 & 39.83 & NC    \\
p3 & NO & 180.55 & 68.84 & 48.83 & NC    \\
p3 & p1 &  68.92 & 58.81 & 46.33 & NC    \\
p3 & p3 &  33.09 & 40.22 & 40.84 & NC    \\\hline
\multicolumn{6}{l}{\textit{\code{diagonalIion}}}\\
NO & NO & 153.16 & 62.86 & 44.92 & NC    \\
NO & p1 &  90.92 & 54.74 & 43.08 & NC    \\
NO & p3 &  32.37 & 38.97 & 39.66 & NC    \\
p3 & NO & 180.42 & 68.71 & 48.67 & NC    \\
p3 & p1 &  68.83 & 58.70 & 46.20 & NC    \\
p3 & p3 &  33.00 & 40.03 & 40.65 & NC    \\\hline
\multicolumn{6}{l}{\textit{\code{JFNK}}}\\
NO & NO & 153.13 & 62.85 & 44.87 & NC    \\
NO & p1 &  90.92 & 54.72 & NC    & NC    \\
NO & p3 &  32.36 & 38.92 & 39.59 & NC    \\
p3 & NO & 180.38 & 68.69 & NC    & NC    \\
p3 & p1 &  68.83 & NC    & NC    & NC    \\
p3 & p3 &  33.00 & NC    & NC    & NC    \\
\bottomrule\end{longtable}\end{center}

\begin{center}\footnotesize
\begin{longtable}{llrrrr}
\caption{Tiempo de reloj de pared del solver 2D [min] a $\Delta t = 10^{-4}$~s
(redondeado a minutos enteros).}\\
\toprule
Vm & Iion & $0.5$ & $0.25$ & $0.125$ & $0.1$ \\
\midrule\endfirsthead\toprule
Vm & Iion & $0.5$ & $0.25$ & $0.125$ & $0.1$ \\
\midrule\endhead
\multicolumn{6}{l}{\textit{\code{Picard}}}\\
NO & NO &  15 &  31 &  83 & 125 \\
NO & p1 &  20 &  51 & 141 & 197 \\
NO & p3 &  23 & 102 & 402 & NC  \\
p3 & NO &  18 &  40 & 112 & NC  \\
p3 & p1 &  20 &  60 & 176 & NC  \\
p3 & p3 &  22 & 105 & 419 & NC  \\\hline
\multicolumn{6}{l}{\textit{\code{diagonalIion}}}\\
NO & NO &  18 &  31 &  81 & NC  \\
NO & p1 &  22 &  52 & 147 & NC  \\
NO & p3 &  29 & 131 & 533 & NC  \\
p3 & NO &  21 &  38 & 107 & NC  \\
p3 & p1 &  23 &  59 & 164 & NC  \\
p3 & p3 &  28 & 133 & 549 & NC  \\\hline
\multicolumn{6}{l}{\textit{\code{JFNK}}}\\
NO & NO &  59 & 155 & 544 & NC  \\
NO & p1 &  82 & 263 & NC  & NC  \\
NO & p3 & 124 & 647 & 3407 & NC \\
p3 & NO &  67 & 169 & NC  & NC  \\
p3 & p1 &  87 & NC  & NC  & NC  \\
p3 & p3 & 131 & NC  & NC  & NC  \\
\bottomrule\end{longtable}\end{center}

\begin{center}\footnotesize
\begin{longtable}{llrrrr}
\caption{Tamaño de conjunto residente pico 2D [MiB] a $\Delta t = 10^{-4}$~s.}\\
\toprule
Vm & Iion & $0.5$ & $0.25$ & $0.125$ & $0.1$ \\
\midrule\endfirsthead\toprule
Vm & Iion & $0.5$ & $0.25$ & $0.125$ & $0.1$ \\
\midrule\endhead
\multicolumn{6}{l}{\textit{\code{Picard}}}\\
NO & NO & 148 & 152 & 174 & 194 \\
NO & p1 & 149 & 157 & 191 & 219 \\
NO & p3 & 155 & 177 & 288 & NC  \\
p3 & NO & 155 & 191 & 346 & NC  \\
p3 & p1 & 156 & 194 & 358 & NC  \\
p3 & p3 & 161 & 218 & 457 & NC  \\\hline
\multicolumn{6}{l}{\textit{\code{diagonalIion}}}\\
NO & NO & 148 & 153 & 176 & NC  \\
NO & p1 & 149 & 158 & 194 & NC  \\
NO & p3 & 155 & 181 & 296 & NC  \\
p3 & NO & 155 & 191 & 347 & NC  \\
p3 & p1 & 156 & 195 & 365 & NC  \\
p3 & p3 & 161 & 218 & 457 & NC  \\\hline
\multicolumn{6}{l}{\textit{\code{JFNK}}}\\
NO & NO & 148 & 153 & 173 & NC  \\
NO & p1 & 149 & 158 & NC  & NC  \\
NO & p3 & 155 & 181 & 299 & NC  \\
p3 & NO & 152 & 187 & NC  & NC  \\
p3 & p1 & 153 & NC  & NC  & NC  \\
p3 & p3 & 160 & NC  & NC  & NC  \\
\bottomrule\end{longtable}\end{center}

\paragraph{Discusión de los resultados 2D.}
A $(\Delta x, \Delta t, $par LRE$)$ fijos las tres estrategias no lineales dan
tiempos de activación que coinciden dentro de $\sim 0.1$--$0.2$~ms; p.\,ej.\ a
$\Delta x = 0.125$~mm con $(\code{NO}, \code{NO})$ los valores son
$45.03/44.92/44.87$~ms para Picard, diagonalIion, JFNK respectivamente. La elección
del acoplamiento no lineal 2D es por lo tanto esencialmente \emph{gratis} desde el
punto de vista de la precisión aquí --- los errores relevantes son espaciales y
temporales, no errores de iteración no lineal. Los tres métodos difieren
marcadamente en costo: JFNK es $\sim 3$--$6\times$ más caro que Picard para el mismo
caso a $\Delta x = 0.5$~mm y crece más rápido que linealmente con el tamaño de malla,
porque cada producto matriz--vector de Krylov realiza una re-evaluación completa del
ODE ($\Delta x = 0.125$, $(\code{NO},\code{p3})$: $402$~min para Picard,
$533$~min para diagonalIion, $3407$~min --- $\sim 2.4$~días --- para JFNK). El RSS
pico está dominado por el stencil LRE y es esencialmente independiente del método; la
columna JFNK es levemente menor para algunas entradas porque la envoltura matrix-free
evita la corrección diagonal y el workspace explícito de BiCGSTAB.

\subsection{Benchmark 3D estilo Niederer}
\label{sec:results:3D}

El barrido 3D sigue la geometría de losa publicada de Niederer
($20\times 3\times 7$~mm, $40\times 6\times 14$ celdas a $\Delta x = 0.5$~mm). Al
momento de escribir, solo los barridos \code{diagonalIion} y \code{Picard} incluyen
corridas 3D, todas a $\Delta x = 0.5$~mm. El barrido \code{JFNK} 3D aún no se ha
corrido.

\begin{center}\footnotesize
\begin{longtable}{llrr}
\caption{Tiempo de activación P8 3D [ms] a $\Delta x = 0.5$~mm,
$\Delta t = 10^{-4}$~s. Las columnas son los métodos no lineales disponibles.}\\
\toprule
Vm & Iion & \code{Picard} & \code{diagonalIion} \\
\midrule\endfirsthead\toprule
Vm & Iion & \code{Picard} & \code{diagonalIion} \\
\midrule\endhead
NO & NO & 136.67 & 136.61 \\
NO & p1 &  33.50 &  33.39 \\
NO & p3 &  30.59 &  30.51 \\
p3 & NO & 149.13 & 149.06 \\
p3 & p1 &  32.44 &  32.36 \\
p3 & p3 &  29.63 & NC     \\
\bottomrule\end{longtable}\end{center}

\begin{center}\footnotesize
\begin{longtable}{llrr}
\caption{Tiempo de reloj de pared del solver 3D [min] a $\Delta x = 0.5$~mm,
$\Delta t = 10^{-4}$~s.}\\
\toprule
Vm & Iion & \code{Picard} & \code{diagonalIion} \\
\midrule\endfirsthead\toprule
Vm & Iion & \code{Picard} & \code{diagonalIion} \\
\midrule\endhead
NO & NO &   80 &   82 \\
NO & p1 &  168 &  229 \\
NO & p3 & 1489 & 2159 \\
p3 & NO &  163 &  143 \\
p3 & p1 &  167 &  269 \\
p3 & p3 & 1487 & NC   \\
\bottomrule\end{longtable}\end{center}

\begin{center}\footnotesize
\begin{longtable}{llrr}
\caption{Tamaño de conjunto residente pico 3D [MiB] a $\Delta x = 0.5$~mm,
$\Delta t = 10^{-4}$~s.}\\
\toprule
Vm & Iion & \code{Picard} & \code{diagonalIion} \\
\midrule\endfirsthead\toprule
Vm & Iion & \code{Picard} & \code{diagonalIion} \\
\midrule\endhead
NO & NO &  156 &  157 \\
NO & p1 &  211 &  211 \\
NO & p3 &  787 &  849 \\
p3 & NO & 1123 & 1121 \\
p3 & p1 & 1164 & 1162 \\
p3 & p3 & 1565 & NC   \\
\bottomrule\end{longtable}\end{center}

\paragraph{Discusión de los resultados 3D.}
En la losa 3D gruesa ($\Delta x = 0.5$~mm), activar la LRE solo en la fuente iónica
cambia el tiempo de activación de $\approx 136.7$~ms (\code{NO}/\code{NO}) a
$\approx 33.5$~ms (\code{NO}/\code{p1}) y $\approx 30.6$~ms (\code{NO}/\code{p3}) ---
el mismo efecto cualitativo de alto-orden-en-la-fuente ya visto en el barrido 2D y en
los reportes Niederer explícito/implícito. El caso totalmente de alto orden
(\code{p3}/\code{p3}) alcanza $29.63$~ms con Picard (la contraparte diagonalIion aún
no se ha completado). Los dos métodos no lineales nuevamente coinciden mejor que
$0.1$~ms para cada entrada donde ambos están disponibles; Picard es generalmente más
rápido, con la brecha más notable en el caso más pesado (\code{NO}/\code{p3}:
$1489$~min para Picard vs.\ $2159$~min para diagonalIion). El RSS pico escala con el
tamaño del stencil LRE y es esencialmente independiente del método. El conjunto de
datos es demasiado ralo en $\Delta x$ para sacar conclusiones cuantitativas sobre
convergencia de malla en 3D; la brecha inmediata a llenar es la entrada faltante
\code{p3}/\code{p3} de diagonalIion y un barrido JFNK 3D a la misma malla.

% ================================================================== %
\appendix
\section{Tabla de dimensiones}
% ================================================================== %

\begin{center}
\begin{tabular}{@{}lll@{}}
\toprule
Campo & \code{dimensionSet} de OpenFOAM & SI \\
\midrule
\code{Vm} & $[1, 2, -3, 0, 0, -1, 0]$ & volt \\
\code{Iion} & \code{dimVoltage/dimTime} & \si{\volt\per\second} \\
\code{externalStimulusCurrent} & \code{dimCurrent/dimVolume} &
    \si{\ampere\per\meter\cubed} \\
\code{conductivity} &
    \code{dimCurrent\^{}2 dimTime\^{}3 / (dimMass dimVolume)} &
    \si{\siemens\per\meter} \\
\code{chi} & \code{dimless/dimLength} & \si{\per\meter} \\
\code{Cm} & \code{dimCurrent dimTime / (dimVoltage dimArea)} &
    \si{\farad\per\meter\squared} \\
\bottomrule
\end{tabular}
\end{center}

\section{Mapa de archivos}

\begin{center}
\begin{tabular}{@{}p{0.55\linewidth}p{0.40\linewidth}@{}}
\toprule
Ruta (relativa a \code{cardiacFoam-pc/}) & Rol \\
\midrule
\code{applications/solvers/.../JfNK.C} & Unidad de traducción principal, incluyendo
   el kernel CellML TNNP embebido, la actualización de states, la evaluación de
   corriente, y la integración local del ODE. \\
\code{applications/solvers/.../createFields.H} & Inicialización de campos
   y lectura de diccionarios. \\
\code{tutorials\_electro/.../run\_convergence.py} & Script driver. \\
\code{constant/cardiacProperties} & $\chi$, $C_{m}$, $\Dten$,
   modelo iónico. \\
\code{constant/timeIntegrationProperties} & Solver de ODE, tolerancias,
   tope de sub-pasos. \\
\code{constant/spatialIntegrationProperties} & Esquema temporal, flags HO,
   opciones de JFNK / lineal / búsqueda de línea. \\
\code{constant/stimulusProtocol} & Regiones y temporización del estímulo. \\
\bottomrule
\end{tabular}
\end{center}

% ================================================================== %
\section{Diferencias vs.\ el hermano PETSc}
% ================================================================== %

Algorítmicamente los dos solvers son idénticos. Las únicas diferencias están en el
back-end de álgebra lineal y unas pocas perillas visibles al usuario:

\begin{center}
\begin{tabular}{@{}p{0.42\linewidth}p{0.26\linewidth}p{0.26\linewidth}@{}}
\toprule
Aspecto & Este solver (Eigen) & Hermano PETSc \\
\midrule
Contenedor de matriz dispersa & \code{Eigen::SparseMatrix} & \code{Mat} de PETSc \\
Solver directo disperso & \code{Eigen::SparseLU} & \code{KSPPREONLY+PCLU} \\
Solver iterativo (rama Picard) & \code{Eigen::BiCGSTAB+ILUT} &
    \code{KSP+PC} (catálogo completo de PETSc) \\
Krylov interno de JFNK & GMRES propio & \code{KSPGMRES} sobre \code{MatShell} \\
Precondicionador dentro de JFNK & ninguno (solo Krylov) & cualquier PC de PETSc \\
Dependencia en tiempo de compilación & Eigen3 header-only &
    PETSc + MPI \\
Soporte de corrida paralela & solo serial & MPI-paralelo de fábrica \\
\bottomrule
\end{tabular}
\end{center}

Todos los parámetros relacionados a la convergencia (\code{nonlinearTolerance},
\code{jfnkEpsilon}, \code{jfnkLineSearch*}, \code{jfnkInitGuess*},
\code{jfnkClampODEInput}, \code{nonlinearAcceptUnconverged},
\code{ODE\_*}) tienen los mismos nombres y significados en ambos solvers, así que la
misma configuración de \code{run\_convergence.py} produce una corrida
algorítmicamente equivalente con cualquiera de los dos back-ends (salvo las
diferencias de rendimiento y de escalado paralelo).

% ================================================================== %
\begin{thebibliography}{9}
% ================================================================== %

\bibitem{ttp2004} K.~H.~W.~J. ten Tusscher, D. Noble, P.~J. Noble, A.~V.
    Panfilov, \emph{A model for human ventricular tissue},
    AJP -- Heart and Circulatory Physiology
    \textbf{286} (2004) H1573--H1589.

\bibitem{niederer2012} S.~A. Niederer et al.,
    \emph{Verification of cardiac tissue electrophysiology simulators
    using an N-version benchmark}, Prog.\ Biophys.\ Mol.\ Biol.
    \textbf{104} (2012) 22--48.

\bibitem{knoll2004} D.~A. Knoll, D.~E. Keyes,
    \emph{Jacobian-Free Newton-Krylov methods: a survey of approaches
    and applications}, J.\ Comp.\ Phys.\ \textbf{193}
    (2004) 357--397.

\bibitem{noceWright} J. Nocedal, S.~J. Wright,
    \emph{Numerical Optimization}, 2nd edition, Springer, 2006.

\bibitem{saadbook} Y. Saad,
    \emph{Iterative Methods for Sparse Linear Systems}, 2nd edition,
    SIAM, 2003.

\bibitem{hairer2} E. Hairer, G. Wanner,
    \emph{Solving Ordinary Differential Equations II}, Springer, 1996.

\end{thebibliography}

\end{document}
